% NOTE: merged into Section II "Background and Related Work" (\input'd right after
% 2-background.tex). The section header and the "Sampling configurable systems"
% paragraph were folded into the merged "How prior work counts effort" paragraph
% there. The remaining paragraphs continue Section II.
\paragraph*{Counting effort versus measuring acquisition cost.}
Beyond the configurable-systems predecessors revisited in \Cref{sec:background},
the same count-as-effort convention recurs across tuning, transfer, and
prediction: ``effort,'' ``efficiency,'' and ``speedup'' are often expressed as
counts, that is, measurements, evaluations, queries, or target-environment
samples, because counts are simple to compare and reuse
\cite{nair2020flash,chen2021mmo,chen2024admmo,ye2025dlisa,gong2023dal,gong2024sempl}.
Counts sound like budgets only if one count unit stands for a comparable amount
of acquisition cost. Otherwise, a cost-effective result can be an accuracy
result wearing a cost label. Even work explicitly framed around cost, such as
cost-efficient sampling \cite{sarkar2015}, still prices each configuration as one
unit, whereas we audit its \emph{measured} cost. This does not invalidate prior accuracy results: a
fixed-$n$ comparison answers a legitimate accuracy question, and only the added
efficiency or budget reading needs measured-cost calibration. Some tuning work
already grounds counts in measured wall-clock budgets
\cite{chen2021multiobj,chen2023aspirations}, while our results show why
per-configuration and per-strategy calibration is still needed.

\paragraph*{Surrogate and heterogeneous costs.}
Replacing a real resource with an easy-to-count surrogate is a recurring
validity concern. In mutation testing, for example, the number of mutants can
mislead relative to execution time \cite{harmanMutationFSE2020}. We ask the
analogous question for configuration sampling, but the mechanism and the
contribution differ from that precedent in three ways. First, the driver is
\emph{strategy bias}: sampling strategies systematically steer toward cheap or
costly regions, so equal-size samples from different strategies spend unequal
measured budgets. There is no analogue of a sampling strategy in mutant
counting. Second, the acquisition cost and the prediction target can be
\emph{different} measured quantities (predicting output size while paying
encoding time), an asymmetry with no mutation-testing counterpart. Third, we do
not merely flag the threat: we quantify its decision impact by re-running the
original studies' own Kruskal--Wallis test under measured cost, and we identify
the conditions (cost heterogeneity, strategy bias, and target--cost
asymmetry) under which the proxy breaks. The surface observation that
configurations have heterogeneous cost is not the contribution. The contribution
is that this changes the set of statistically recommendable samplers and the
efficient sets.

\paragraph*{Context, learners, and optimization.}
Prior work shows that sampling conclusions can vary across inputs, workloads,
and target properties \cite{pereira2020,ye2026domland}, and we use these dimensions
to test cost-proxy validity. Learner choice is a separate modeling issue: we
keep literature-aligned learners for replications and use stronger sklearn
baselines as robustness checks or primary learners for newly framed families
\cite{breiman2001random,valov2015,amand2019,bao2018,sklearn_api}. Active
learning and multi-objective optimization offer candidate solutions, but they
inherit the same proxy question through their cost term
\cite{sarkar2015,siegmund2015FSE,li2020pareto}. We therefore treat measured
cost as a prerequisite for accuracy--cost claims, not as a new optimizer. A
complementary strand instead measures model-building or adaptation cost rather
than label acquisition: training time \cite{gong2022encoding}, online-drift
overhead \cite{xiang2026dhda}, and deep-learning costs \cite{gong2024dcpl}.
Extending our audit to those axes is future work.
